\documentclass{handout}

\assignment{Homework 7}
\title{Electric field}

\begin{document}
\maketitle
\practiceinstructions

\section{Charge and Coulomb's law}

% No old question covered Coulomb's law on its own.

\section{Reading a field}

\begin{questions}
    \question \textbf{Charge on a die}: You come across the following die, with its faces numbered 1-6 in the usual way (opposite faces add to 7). Upon measuring the electric field, you find that at all locations the electric field points to the right (in the $+y$ direction) with varying strengths. On the left of the die it has strength $|\vec{E}|=E_{3}$, on the right $|\vec{E}|=E_{4}$, and on all other faces $|\vec{E}|=E_{other}$. The die is 1.5 cm on each side.
    \begin{figure}[h!]
        \centering
        \includegraphics[height=2in]{Images/GaussDice.pdf}
    \end{figure}
    \begin{parts}
        \part Write an expression ($\pm \hat{x}, \pm \hat{y}$, etc.) for the outward pointing normal vector of each face (1-6).
        \part Write expressions for the dot product $(\vec{E} \cdot \hat{n})$ for each face (hint: some faces should be zero)
        \part Write expressions for the electric flux over each face.
        \part If $E_3=1.2\times 10^6$ N/C, $E_{other}=1.8\times 10^6$ N/C, and $E_{other}=2.5\times 10^6$ N/C, how much charge is inside the cube? (Remember $\Phi_E=Q_{in}/\epsilon_0$ where $\epsilon_0= 8.854\times 10^{-12}$ F/m.
    \end{parts}
    \question \textbf{Potential fields}: For the following potential field, mark where there are regions of positive and negative charge. Explain your reasoning in terms of gradients, the electric field, and divergence.
    \begin{figure}[h!]
        \centering
        \includegraphics[height=4in]{Images/Peaks.png}
    \end{figure}
\end{questions}

\end{document}
